% ------------------------------------------------------------------------------
% SETUP ------------------------------------------------------------------------
% ------------------------------------------------------------------------------

\documentclass[12pt]{article}
\usepackage[a4paper, width = 175mm, top = 20mm, bottom = 20mm, 
bindingoffset = 0mm]{geometry}
\usepackage[utf8]{inputenc}
\usepackage{url}
\usepackage{ragged2e}
\usepackage{array}
\usepackage{tabularx}
\usepackage{float}
\usepackage{xcolor, listings}
\usepackage{textcomp}
\usepackage{fancyhdr}
\newcommand{\changefont}{%
    \fontsize{8}{11}\selectfont
}

\usepackage[hang, flushmargin, bottom]{footmisc} 
\usepackage{hyperref}
\hypersetup{
  colorlinks = true,
  linkcolor = .,
  urlcolor = .,
  citecolor = .,
  bookmarks = true}

\pagestyle{fancy}
\fancyhead{}
\fancyhead[R]{\changefont{Topic-specific sentiment analysis for tweets by 
German MPs}}
\fancyfoot{}
\fancyfoot[R]{\thepage}
\setlength{\headheight}{14.5pt}
\setlength{\parindent}{0pt}
\interfootnotelinepenalty = 10000

% ------------------------------------------------------------------------------
% MAIN -------------------------------------------------------------------------
% ------------------------------------------------------------------------------

\usepackage{/usr/share/R/share/texmf/tex/latex/Sweave}
\begin{document}
\Sconcordance{concordance:einladung.tex:einladung.Rnw:%
1 43 1 1 0 94 1}


% FRONT PAGE -------------------------------------------------------------------
 
\begin{titlepage}
\begin{center}
    
\large
Statistical Consulting
    
\vspace{0.3cm}
      
\rule{\textwidth}{1pt}
% \large
\textbf{Topic-specific sentiment analysis for tweets by German MPs}
\rule{\textwidth}{1pt}
   
\vspace{0.3cm}
      
\normalsize
Department of Statistics \\
Ludwig-Maximilians-Universität München 

\vspace{0.3cm}

\includegraphics[width = 0.2\textwidth]{figures/sigillum.png}

\vspace{0.3cm}

\rule{\textwidth}{1pt}
\small
% \raggedright
New data sources offer vast potential for studying public opinion in the 
political environment.
Frequently arising research questions concern the nature of topics discussed 
and the sentiment expressed toward them.
We perform sentiment analysis for tweets by German MPs issued since the last 
federal election in September 2017, scraped from the Web, with special regard to 
topical context.
The contribution of this work is twofold.
First, we explore how such analysis can be conducted with either standard 
machine learning approaches for tabular data or more complex, BERT-based deep 
learning solutions.
Our results suggest that casting the problem as a standard machine 
learning task achieves acceptable 
predictive performance and requires extensive feature engineering.
BERT-based models clearly outperform this simplistic approach while avoiding 
the need for hand-crafted inputs.
The ability to draw on knowledge acquired during the pre-training phase possibly 
supports the Transformers in solving the difficult task of classifying 
relatively few observations in a sparsely populated feature space.
Post-training on the large pool of unlabeled data typically available in such 
settings is also helpful.
We further note that topic modeling does not aid classification in this 
application, in fact complicating it, but attribute this largely to the fact 
that extracting topics from this kind of data is challenging even for the 
human annotator.
Second, we provide teaching material for instructing fellow researchers in the 
applied techniques.
With this collection, which will remain publicly available, we hope to 
facilitate the entrance to NLP for practitioners from other disciplines.
\rule{\textwidth}{1pt}

% \center

\vspace{0.3cm}

\begin{tabular}{rl}
  Speakers & Asmik Nalmpatian \& Lisa Wimmer \\
  & \\
  Project partner & Prof. Dr. Paul Thurner \\
  & \textit{Department of Political Science} \\
  & \\
  Supervisors & Prof. Dr. Christian Heumann\\
  & Matthias Aßenmacher \\
  & \textit{Department of Statistics} \\
  & \\
  Date & Monday, July 12\textsuperscript{th}, 2021, 16:00 s.t. \\
  & \\
  Meeting & ID 918 5472 5495, passcode 413775
\end{tabular}

\vspace{0.3cm}

\rule{\textwidth}{1pt}

\vspace{0.3cm}

\url{https://lmu-munich.zoom.us/j/91854725495?pwd=NHVLemt4ZzZHNk8vSUpHWHBZamZjUT09}

\end{center}
\end{titlepage}

\end{document}
